\documentclass[conference]{IEEEtran}
\IEEEoverridecommandlockouts

\usepackage{amsmath,amssymb,amsthm}
\usepackage{algorithm}
\usepackage{algpseudocode}
\usepackage{listings}
\usepackage{xcolor}
\usepackage{booktabs}
\usepackage{graphicx}
\usepackage{hyperref}
\usepackage{cite}
\usepackage{enumitem}

\newtheorem{theorem}{Theorem}
\newtheorem{definition}{Definition}
\newtheorem{invariant}{Invariant}
\newtheorem{proposition}{Proposition}

\lstdefinelanguage{JavaScript}{
  keywords={function,async,await,if,else,const,let,var,return,export,import,from},
  keywordstyle=\color{blue}\bfseries,
  comment=[l]{//},
  commentstyle=\color{gray}\itshape,
  stringstyle=\color{red},
  morestring=[b]',
  morestring=[b]"
}
\lstset{
  language=JavaScript,
  basicstyle=\ttfamily\footnotesize,
  breaklines=true,
  frame=single,
  numbers=none,
  columns=flexible
}

\begin{document}

\title{Static Reactive Linking for Plain TypeScript/JSX:\\
Cross-Module Invalidation without Runtime Read Tracking}

\author{
\IEEEauthorblockN{Amihere Theophilus Junior}
\IEEEauthorblockA{
Department of Computer Science and Engineering\\
University of Mines and Technology (UMaT)\\
Tarkwa, Ghana\\
Email: snrraptopack@gmail.com}
}

\maketitle

\begin{abstract}
Every fine-grained reactive framework must eventually answer the same
question after a mutation: which computations are now stale, and in what
order should they be brought up to date? Signal-based systems answer this
dynamically, by recording a subscription every time a reactive value is read.
Static approaches predate signals by decades: IceDust derives incremental
maintenance strategies from static object paths, PixieDust carries that
analysis into browser views, and Svelte's original compiler instruments plain
assignments and orders reactive declarations before the program ever runs.
What remains less explored is how far static analysis can be carried through
\emph{ordinary}, mutable TypeScript/JSX code -- code whose state, mutator
functions, derived values, components, and effects are not confined to one
file but scattered across an ES-module graph the way real applications are
actually written. This paper studies that question directly and presents
\emph{static reactive linking}, implemented in the \texttt{memoized-dom}
prototype. A binding-aware analysis extracts, for every computation, its read
set, and for every function, a summary of its write effects. A fixed-point
module linker gives every imported binding a canonical, alias-independent
identity tied to its defining module, and propagates exact, receiver-bounded,
and parameter-relative effect summaries through chains of imported helper
functions. Each compiled module contributes a fragment to a shared dependency
index mapping canonical state paths to structural UI-entity patterns; at
runtime, a mutation consults this compiler-generated index and a live entity
registry, while an ordinary reactive \emph{read} never allocates or updates a
subscription. We give a module-aware formal model of this reactive identity
and invalidation scheme, and prove dirty-set coverage conditional on five
explicit soundness obligations -- read coverage, write coverage, summary
soundness, commit instrumentation, and pattern--instance consistency --
together with a closed-world linking assumption. We state these assumptions
explicitly rather than treating them as implicit. We are equally explicit
about the cost of this design: static
linking trades read-time dependency discovery for conservative branch unions,
structural wildcard expansion, and whole-subtree fallbacks wherever an effect
cannot be bounded. In a purposive corpus of 88 write-effect sites, linking
reduced the unbounded tier from 45.45\% to 0.00\%; across 25 concrete oracle
executions, static selection achieved 100\% observed recall and 80.0\%
precision. A rendering-free mechanism benchmark found no uniform routing
advantage: the static median was lower in one of six topologies and higher in
the other five, with lifecycle cost depending strongly on wildcard density.
\end{abstract}

\begin{IEEEkeywords}
reactive programming, static analysis, module linking, incremental user
interfaces, effect summaries, TypeScript, JSX
\end{IEEEkeywords}

\section{Introduction}

Consider a small, entirely ordinary refactor. A developer pulls a piece of UI
state -- a selected row ID, say -- out of the component that first needed it
and into a shared module, so that three unrelated components can read it and
one shared helper function can update it. In a signal-based framework this
refactor is invisible to the reactivity system: wherever the value is read,
a subscription is recorded at that moment, regardless of which file the read
happens to live in, or how many function calls separate the read from the
component that renders it. The dependency graph is discovered lazily, one
execution at a time, and it does not care about module boundaries because it
is built \emph{after} the module system has already resolved everything into
running code.

A compiler that wants to determine the same dependency relationship
\emph{before} the program runs does not get this for free. It has to answer,
statically, exactly the questions the module system answers dynamically:
which exported binding does this import alias refer to, does this exported
mutator function actually write the state it appears to write, and does a
helper three calls deep eventually touch a property on an object handed to it
by its caller? Static dataflow and incremental-computation research answered
versions of this question long before signals existed. Lustre compiles
synchronous equations into bounded-memory sequential code \cite{lustre};
Incremental $\lambda$-Calculus statically differentiates pure programs into
their own incremental updaters \cite{incremental-lambda}; IceDust inverts
static object-graph paths into dependency-ordered maintenance strategies
\cite{icedust}; PixieDust carries that same path analysis into an actual
browser UI \cite{pixiedust}; and Svelte's compiler has, since version 3,
instrumented ordinary variable assignments and statically ordered its
reactive declarations \cite{svelte3,svelte-legacy}. These systems do not
directly target the refactor above -- a reactive value moved
across an ES-module boundary, read through an import alias, and written
through an imported helper function -- in the setting of ordinary, mutable
TypeScript and JSX rather than a purpose-built DSL or a single-file component
model. That is the gap this paper works in, and it is deliberately narrower
than ``static reactivity is possible'':

\begin{quote}
Can a compiler preserve statically inferred reactive identity and mutation
effects through an ordinary ES-module graph, then route writes to live UI
instances without discovering dependencies at read time?
\end{quote}

The question matters because application state rarely stays where it was
first declared. It gets exported from dedicated state modules, mutated
through layers of imported helper functions, transformed into module-owned
derived values, and consumed by components and effects declared in files that
may never import one another directly. ES modules complicate this further
with aliasing: two files can export a binding under the same local name, and
one exported object can be re-imported under several different names, so a
naive syntactic match on identifier text is not enough to recover a single
reactive identity.

This paper makes five contributions:

\begin{enumerate}
  \item A module-aware model in which a reactive location is identified by
  its defining module, binding, and static member path -- independent of
  whatever alias a particular importer happens to use.
  \item A fixed-point linker that propagates binding-aware read effects,
  exact and bounded writes, and parameter-relative effects through named
  ES-module imports and chains of helper functions.
  \item A compiler/runtime design that maps those canonical locations onto
  exact and structural live UI entities -- including module-owned derived
  values and module-owned effects -- without tracking reads at runtime.
  \item A conditional invalidation-coverage theorem that names the precise
  soundness obligations the compiler must discharge, rather than treating
  static-analysis completeness as an unexamined given.
  \item An artifact-backed evaluation that separates coverage, selection
  precision, runtime cost, and the specific value and cost contributed by
  cross-module linking.
\end{enumerate}

We use \texttt{memoized-dom} as the reference implementation throughout. To
be direct about what is \emph{not} claimed: we do not claim it accepts full
ECMAScript, that it stores no runtime structures, or that it has already been
shown faster than signal-based systems. Each of these is either false as
stated or an open empirical question, and Sections~\ref{sec:runtime}
through~\ref{sec:limits} say so explicitly.

\section{Related Work}
\label{sec:related}

Positioning this work correctly means resisting a tempting but incorrect
shortcut: treating ``static reactivity'' as a single, already-claimed idea
and ``generated DOM updates'' as another. Neither is a sufficient novelty
claim on its own -- both have a long and specific prior history, and the
honest way to describe this paper's place in that history is to walk through
each precedent and say exactly what it already establishes and exactly where
it stops.

\subsection{Static Dataflow and Incremental Transformation}

Lustre is the oldest relevant precedent and the furthest from UI programming:
a synchronous dataflow language for reactive control and monitoring systems,
in which equations and explicit clocks admit compilation into bounded-memory
sequential or automaton implementations \cite{lustre}. It establishes, as
far back as 1991, that a statically structured reactive network can be
compiled rather than interpreted at runtime -- but it says nothing about
inferring such a network from mutable, imperative ES-module code, and nothing
about DOM instance lifecycles, which did not yet exist as a problem.

Incremental $\lambda$-Calculus (ILC) moves the idea into a higher-order
functional setting: a pure program $f$ is statically differentiated into an
incremental function mapping input changes to output changes, automatically,
without any runtime dependency graph, and with a machine-checked correctness
development for a family of typed calculi \cite{incremental-lambda}. This is
a broad and rigorous precedent for the claim that static transformation can
replace runtime tracing -- but ILC computes general output \emph{deltas} for
pure terms, whereas the system studied here attributes imperative
\emph{writes} to the computations they affect and typically re-evaluates a
guarded expression rather than deriving a closed-form delta. ILC is a
theoretical ancestor, not an implementation-equivalent system.

Self-adjusting computation and its systems descendant Adapton sit at the
opposite end of the runtime-cost spectrum from both of the above: they
construct or maintain execution-derived dependency traces \emph{at run time},
supporting far more general computation than the restricted static surface
studied here, at the cost of the very runtime bookkeeping this paper tries to
avoid \cite{acar-sac,adapton}.

\subsection{Path-Based Static Dependency Analysis}

IceDust is the closest formal-analysis ancestor to this work, and it is worth
describing in some detail rather than summarizing in a single sentence. It
extracts static paths from derived-attribute expressions in a declarative,
persistent object-graph DSL, inverts those paths into a dataflow relation,
computes strongly connected components and a topological order over the
resulting graph, and generates one of three concrete maintenance strategies
-- Calculate-on-Read, Calculate-on-Change, or Calculate-Eventually -- per
derived attribute \cite{icedust}. Every one of path extraction, dependency
inversion, generated incremental maintenance, and topological ordering is
already established by IceDust. What IceDust does not do is address JSX DOM
entities, ordinary imperative TypeScript mutation, ES-module alias
resolution, or component placement across files, because its derived
expressions are read-only by construction and its source language is a
purpose-built DSL rather than a general-purpose imperative language.

PixieDust takes the next step and is the nearest UI precedent of all: it
extends IceDust's path analysis with functions and views, statically
over-approximates model/view dependencies, inverts them, and uses the result
to incrementally update a React-rendered browser UI, reporting performance
comparable to MobX in its own evaluation \cite{pixiedust}. Because it is this
close, it deserves to be discussed on its own terms rather than folded into a
generic ``incremental computation'' paragraph. Two things distinguish this
paper's approach from PixieDust. First, PixieDust's runtime is not empty --
its own formal semantics describes per-field component subscriptions, a
component metadata store, dirty flags, a re-render queue, and a cached
virtual DOM -- which is an important reminder, discussed further in
Section~\ref{sec:runtime}, that static dependency \emph{analysis} does not by
itself imply the absence of runtime dependency \emph{structures}. Second,
PixieDust's source language is an explicit model/view/action DSL with
declared inverse relationships between entities, whereas this work analyzes
ordinary TypeScript/JSX, infers imperative writes rather than requiring
declared ones, and links reactive identity across ES-module boundaries that
PixieDust's single-model setting does not need to cross.

Reactive variables integrate signal-like values directly into an imperative
language and use lexical restrictions to guarantee that the resulting
dependency graph is acyclic and can be topologically ordered
\cite{reactive-variables}. This is a close precedent at the language-design
level, though it is not centered on DOM compilation or on routing writes to
structural instances placed across separately compiled files.

\subsection{Industrial UI Compilers}

Of all the systems compared here, legacy Svelte (versions 3 and 4) is the
one closest to production reality, and the one whose documented limitations
are most instructive. Its compiler instruments ordinary assignments directly
-- \texttt{count += 1} becomes an invalidation call -- and determines the
dependencies of \texttt{\$:} reactive statements purely at compile time,
ordering them topologically before the program runs \cite{svelte3,svelte-legacy}.
Svelte's own documentation is candid about where this breaks down: a
dependency hidden behind a function call is invisible to the analysis unless
it appears lexically inside the reactive statement itself, which is precisely
the class of case a linker across ES-module helper chains is built to
recover. The other half of the comparison is just as telling -- shared state
that lives outside a single component normally falls back to Svelte's
\texttt{store} primitive and its own runtime subscription contract, which is
a tacit admission that compile-time analysis, in that design, does not extend
cleanly across module boundaries. This paper's linker is, in effect, an
attempt to push exactly that boundary further out.

Svelte 5, Solid, and Vue represent a different and more recent design choice:
they deliberately determine dependencies \emph{dynamically}, at evaluation
time, using signals, refs, or reactive proxies under the hood
\cite{svelte-runes,solid-reactivity,vue-reactivity}. Solid registers the
currently executing computation as an observer the instant a signal getter
runs; Vue tracks proxy and ref reads in a nested target/key/effect structure
and replays stored effects on write; Svelte 5's \texttt{\$derived} and
\texttt{\$effect} resolve their dependencies when evaluated, using signals
internally. The appeal of this design is real: dependency discovery is
path-sensitive by construction, so a dependency hidden behind an arbitrary
function call is simply observed if it is actually read on a given execution,
with no special-casing required. The cost is a runtime dependency
representation that must be built and kept consistent as the program runs.
Static linking accepts the opposite trade: a conservative, closed-world
over-approximation computed once at compile time, in exchange for removing
dependency \emph{discovery} from the read path -- and, as
Section~\ref{sec:limits} makes clear, an explicit fallback surface wherever
that approximation cannot be computed precisely.

\subsection{Ahead-of-Time Reactive Scheduling}

Outside the browser, Oeyen et al. present Remus, a low-level reactive virtual
machine influenced by Haai, in which reactor graphs can be topologically
pre-scheduled ahead of time -- avoiding a dynamically maintained priority
queue -- while dedicated deployment commands still permit controlled dynamic
structural change \cite{remus}. Remus targets glitch freedom and bounded
overhead on constrained embedded hardware, working from an explicit reactor
description rather than inferring one from imperative source code; its
comparison here is closer in spirit (ahead-of-time scheduling as a
performance strategy) than in mechanism.

\subsection{Where This Leaves the Contribution}

Table~\ref{tab:positioning} lays the whole comparison out along the axis that
matters most for this paper: not whether a system is ``static'' or
``dynamic'' in the abstract, but specifically \emph{when} the
state-to-computation dependency relationship is determined, and \emph{what},
if anything, a running program must still keep in memory to route a write.
Read this way, the contribution is not a new point on the static/dynamic
axis -- IceDust and PixieDust already occupy the static end of it for their
respective source languages -- but an attempt to carry that same static
determination through an ordinary, multi-file, imperative TypeScript/JSX
program, which is a source-language problem none of these systems needed to
solve for their own settings.

\begin{table*}[t]
\centering
\caption{Positioning by dependency determination and program model}
\label{tab:positioning}
\begin{tabular}{@{}p{2.4cm}p{3.0cm}p{3.0cm}p{3.5cm}p{3.5cm}@{}}
\toprule
\textbf{System} & \textbf{Source model} & \textbf{Dependency determination} & \textbf{Cross-unit abstraction} & \textbf{Run-time dependency state} \\
\midrule
Lustre & synchronous streams/equations & static & dataflow nodes/modules & compiled automaton/stream state \\
ILC & pure typed $\lambda$ terms & static differentiation & higher-order composition & no dependency graph required \\
IceDust & object-graph DSL & static paths and inversion & whole data model & strategy-specific cache/scheduler \\
PixieDust & model/view/action UI DSL & static paths and inversion & functions and entity views & field subscriptions, component store, queue, VDOM cache \\
Svelte 3/4 & component language & static assignments and \texttt{\$:} & stores for shared state & component dirty state; store subscriptions \\
Solid/Svelte 5/Vue & signals/runes/refs/proxies & dynamic reads & reactive values cross functions/modules & source--observer graph/maps \\
Remus/Haai & reactors & static pre-scheduling plus deployments & reactor composition & VM deployment state \\
\texttt{memoized-dom} & TS/JSX subset & static reads/writes and linking & canonical ES-module locations and summaries & static index, wildcard expansions, entity registry, dirty queue \\
\bottomrule
\end{tabular}
\end{table*}

\section{Source and Compilation Model}
\label{sec:source}

\subsection{Analyzed Source Subset}

The prototype accepts TypeScript/JSX broadly, but it only assigns reactive
meaning to a piece of code where its effects can be bounded statically --
everything else falls through to the conservative handling described in
Section~\ref{sec:limits} rather than being silently ignored. Recognized state
includes module-owned \texttt{let}/\texttt{var} bindings, mutable
object/array roots, component-local state, props, and expressions the
compiler recognizes as derivations of that state. Critically, reactivity is
attached to ordinary JavaScript mutation semantics rather than to a required
API: a state root is reactive whether it is updated by direct assignment
(\texttt{count++}, \texttt{store.selectedId = id}) or by one of the standard
mutating array/object operations (\texttt{push}, \texttt{splice},
\texttt{sort}, \texttt{delete}, \texttt{Object.assign}, and similarly for
\texttt{Map}/\texttt{Set} mutators). The compiler recognizes a fixed table of
such mutating built-ins and attributes their effect to the receiver's
canonical root path exactly as it would an assignment -- so
\texttt{store.todos.push(item)} is analyzed as a write to
\texttt{state.ts\#store.todos}, with no wrapper, decorator, or explicit
signal call required at the call site. This is what makes the source subset
``ordinary'' in the sense claimed throughout this paper: a developer writes
plain mutable state and plain mutating operations, and reactivity is a
property the compiler discovers, not an API the developer opts into. The
analysis follows lexical binding identity, static member paths, local
aliases, and summaries computed for every function visible to it. Dynamic
reflection, computed member mutation through a non-static key, and behavior
hidden inside unresolved external code sit outside this precise subset by
design, not by oversight.

\subsection{Canonical Module State}

Let $M$ be the finite set of modules in one connected compilation graph. The
first problem the linker has to solve is the aliasing problem sketched in the
introduction: the same underlying state must resolve to the same identity no
matter which file reads it or what local name it was imported under. A
canonical reactive location makes this precise:

\begin{definition}[Canonical reactive location]
\label{def:key}
For defining module $m\in M$, declared root binding $b$, and possibly empty
static member path $\pi$,
\[
  \operatorname{key}(m,b,\pi)
  = \operatorname{moduleId}(m)\mathbin{\#}b[.\pi].
\]
\end{definition}

Import resolution maps a local alias to the key of its defining export, which
is what prevents both false separation of a value imported under two names,
and accidental collision between two unrelated exports that happen to share a
name.

\begin{lstlisting}[caption={Cross-module state, derivation, and view}]
// state.ts
export let count = 0;
export function increment() { count++; }

// derived.ts
import { count } from './state';
export const parity = count % 2 ? 'odd' : 'even';

// App.tsx
import { increment } from './state';
import { parity } from './derived';
export function App() {
  return <button onClick={increment}>{parity}</button>;
}
\end{lstlisting}

This three-file example is the refactor from the introduction made concrete.
\texttt{increment} writes \texttt{state.ts\#count}; a singleton computation in
\texttt{derived.ts} reads that same canonical key and conditionally commits
\texttt{derived.ts\#parity}; and the component declared in \texttt{App.tsx}
reads that second key in turn. Three files, two hops of derivation, and one
canonical identity threading all of it together -- which is exactly the chain
a purely lexical, single-file analysis cannot see.

The same canonical-identity machinery applies unchanged when the mutation is
a built-in array or object operation rather than an assignment:

\begin{lstlisting}[caption={Mutating-method write via an imported helper}]
// state.ts
export const store = { todos: [] as string[] };
export function addTodo(text: string) {
  store.todos.push(text); // mutating method, not '='
}

// TodoList.tsx
import { store } from './state';
export function TodoList() {
  return <ul>{store.todos.map(t => <li>{t}</li>)}</ul>;
}
\end{lstlisting}

Here \texttt{addTodo} is linked as a write to
\texttt{state.ts\#store.todos} even though no assignment operator appears in
its body -- the compiler's write-extraction pass recognizes
\texttt{.push} as a mutating operation on a statically resolved receiver and
attributes the effect to that receiver's canonical path, then propagates that
summary through \texttt{addTodo} to any importer that calls it, precisely as
it would for an exact assignment.

\subsection{Function Effect Summaries}
\label{sec:tiers-ref}

For each analyzed function $f$, the compiler computes

\[
S(f)=\langle R_f,W_f,B_f,Q_f,U_f\rangle,
\]

where $R_f$ is the read set, $W_f$ exact writes, $B_f$ effects bounded to
known receivers/state roots -- which includes both direct property
assignment and calls to the compiler's recognized table of mutating
built-ins (\texttt{push}, \texttt{pop}, \texttt{splice}, \texttt{sort},
\texttt{reverse}, \texttt{delete}, \texttt{Object.assign}, and the
corresponding \texttt{Map}/\texttt{Set} mutators) on a statically resolved
receiver -- $Q_f$ paths relative to formal parameters, and $U_f$ an
unbounded-effect flag. At a call site, parameter-relative paths are
substituted with the statically known origins of the actual arguments, so
that a helper written once and called from several places contributes a
different concrete write-set at each call, without needing to be re-analyzed
from scratch.

The linker itself proceeds as a worklist algorithm: it first discovers
exports and imports across the module graph, then repeatedly re-analyzes
modules against the current summaries. If a module's exported summary
changes as a result, every module that imports it is requeued for
re-analysis. Because the abstract facts being propagated are finite sets over
finite module syntax plus a single Boolean unbounded flag, this process is
expected to converge, and the implementation additionally enforces a
diagnostic iteration bound as a safety net. A full monotonicity proof for
every individual analysis transfer function remains future work, and we say
so plainly rather than asserting termination as already established.

\subsection{Structural Computation Identity}

Computations are DOM update closures, conditional/list regions, component
updates, module computeds, or effects. Because a single component
declaration can be instantiated from more than one call site, a computation
declaration does not always correspond to a single runtime placement -- the
compiler therefore emits a \emph{set} of patterns describing every place a
given declaration might live:

\[
p ::= s_1/\dots/s_n
  \mid s_1/\dots/\mathrm{Row}[*]
  \mid s_1/\dots/s_n/*.
\]

Exact patterns name singleton placements; row wildcards denote keyed
instances within a list; descendant patterns conservatively cover an entire
structural subtree when nothing more precise can be established. Component
placement itself is linked through the cross-file component call graph, so
the pattern set emitted for a component declared in one module is often
derived from callers declared in others -- placement, like state, does not
respect file boundaries, and the analysis has to follow it across them.

\section{Static Dependency Index and Runtime}
\label{sec:runtime}

\subsection{Index Synthesis}

Let $C$ be the computations, $K$ canonical locations, and $R(c)\subseteq K$
the statically inferred reads of $c$. Let $\operatorname{Pat}(c)$ be its set
of possible placement patterns.

\begin{definition}[Access index]
\label{def:index}
\[
AT(k)=\bigcup_{c\in C:k\in R(c)}\operatorname{Pat}(c).
\]
\end{definition}

Each compiled module installs its own fragment of $AT$, and the runtime
merges fragments by canonical key as modules initialize. Exact readers are
retained as concrete entity IDs; wildcard patterns are compiled once and then
incrementally maintained as live entities register and unregister, rather
than being re-matched from scratch on every write. Member paths use
segment-aware prefix matching in both directions: a write to a root reaches
readers of its known members, and a write to a member reaches readers of its
root.

It is worth being precise about what this design does and does not eliminate,
because it is easy to overstate. The installed index, the wildcard match
sets, a live entity registry/tree, dirty queues, and caches are all runtime
state, and all of them persist for the life of the application. What is
genuinely absent is dependency \emph{discovery} triggered by a reactive read
-- an ordinary read never allocates, never registers a subscription, and
never touches this runtime state at all. That distinction is why we describe
this design as working ``without runtime read tracking'' rather than
``without a runtime graph''; the latter phrase would not survive contact with
Definition~\ref{def:index}.

\subsection{Mutation Commit}

The compiler instruments an analyzed mutation scope with a hoisted constant
write set:

\begin{lstlisting}[caption={Generated commit for an imperative mutation scope}]
const WRITES_0 = ['./state.ts#count'];
function increment() {
  count++;
  commitWrites(WRITES_0);
}
\end{lstlisting}

\begin{algorithm}[h]
\caption{Static write routing}
\label{alg:commit}
\begin{algorithmic}[1]
\Procedure{CommitWrites}{$W$}
  \If{some $k\in W$ is unbounded}
    \State mark the application root subtree dirty
  \Else
    \State $P\gets\bigcup_{k\in W} AT(k)$, including path-prefix matches
    \State $D\gets\{\ell\in L\mid\exists p\in P:p\sim id(\ell)\}$
    \For{each $\ell\in D$}
      \State mark $\ell$ dirty
    \EndFor
  \EndIf
\EndProcedure
\end{algorithmic}
\end{algorithm}

The dirty set deduplicates marks as it is built. A commit pass processes
ordinary render/computed work before effects, and orders a given batch by
structural depth, so that a computation is not re-evaluated while one of its
own statically known upstream dependents is still settling. Module computeds
are assigned a distinguished early depth so they resolve before the
components that read them; writes produced while a batch is already running
are drained in subsequent passes rather than interleaved unpredictably with
the current one. We want to flag directly that this is not yet a proof of
general glitch freedom over arbitrary cyclic derivations -- cycles are
currently caught by a bounded-pass diagnostic rather than ruled out
structurally, and a genuine topological-ordering proof over the full
computation graph is left as future work.

\subsection{Guards and Conservative Selection}

Because static reads are unions over every syntactically reachable branch of
a computation, a selected computation can turn out to be a false positive
relative to any one concrete execution -- the analysis may correctly note
that a branch \emph{could} read a value without knowing that this particular
run took the other branch. Generated update slots absorb this by comparing
the next value against a cached value before applying any output operation:

\begin{lstlisting}[caption={Guarded output slot}]
if ($slot !== ($next = expression())) {
  $slot = $next;
  applyOutput($next);
}
\end{lstlisting}

This suppresses the externally visible effect of an unchanged output, but it
does not recover the computation cost already paid to evaluate
\texttt{expression()} in the first place -- an important distinction for
Section~\ref{sec:evaluation}, since ``selected but suppressed'' and ``never
selected'' are different costs that a fair evaluation has to measure
separately, not collapse into a single number.

\section{Module-Aware Invalidation Model}
\label{sec:formal}

Having built up the machinery informally, we now state precisely what it
guarantees -- and, just as importantly, exactly what it assumes in order to
guarantee it. Let $L$ be the currently live computation instances. Each
instance $\ell$ is an instantiation of declaration
$\operatorname{decl}(\ell)\in C$ and has a run-time ID $id(\ell)$.

\begin{invariant}[Pattern--instance consistency]
\label{inv:pattern}
For every $\ell\in L$, at least one
$p\in\operatorname{Pat}(\operatorname{decl}(\ell))$ satisfies
$p\sim id(\ell)$.
\end{invariant}

The existential here is not a weakening for convenience; it reflects a real
fact about the system. A reusable component declaration can have several
alternative caller-derived patterns emitted for it, since it may be mounted
from more than one call site -- but any \emph{particular} live instance of
that declaration occupies exactly one placement, and Invariant~\ref{inv:pattern}
only needs one of its patterns to actually match.

For executed scope $\sigma$, let $W_\sigma$ be the write-set used by its
generated commit and define

\[
D_\sigma=\{\ell\in L\mid
  \exists k\in W_\sigma, p\in AT(k):p\sim id(\ell)\}.
\]

\begin{theorem}[Conditional dirty-set coverage]
\label{thm:coverage}
Assume (i) read coverage: every reactive location that may influence $c$ is in
$R(c)$; (ii) write coverage: every relevant mutation performed by $\sigma$ is
in $W_\sigma$ or selects the root fallback; (iii) linked summaries
over-approximate their implementations after argument substitution; (iv) the
scope executes its generated commit at the required boundary; (v)
Invariant~\ref{inv:pattern}; and (vi) the linked module graph is closed over
imports used for precise analysis. Then for every $\ell\in L$,
\[
(R(\operatorname{decl}(\ell))\cap W_\sigma\neq\varnothing)
\implies \ell\in D_\sigma.
\]
For an unbounded effect, marking the root subtree covers every live application
entity.
\end{theorem}

\begin{proof}
Choose $k\in R(\operatorname{decl}(\ell))\cap W_\sigma$. By
Definition~\ref{def:index},
$\operatorname{Pat}(\operatorname{decl}(\ell))\subseteq AT(k)$. By
Invariant~\ref{inv:pattern}, some pattern $p$ in that set matches $id(\ell)$.
Since $k\in W_\sigma$, the definition of $D_\sigma$ includes $\ell$. If the
effect is unbounded, the root-subtree operation selects all live descendants by
construction.
\end{proof}

It is worth dwelling on what this proof does and does not show, because the
gap between the two is exactly where a real compiler can fail silently. The
theorem proves that \emph{if} assumptions (i)--(vi) hold, the routing
mechanics of Section~\ref{sec:runtime} correctly deliver every affected
computation into the dirty set; it is a statement about the algorithm, not
about any particular JavaScript program fed into it. Whether a given analysis
run actually \emph{satisfies} those six assumptions is a separate, harder
question, and one that does not have a uniform yes for arbitrary JavaScript.
Exception exits, opaque callbacks retained by unanalyzed code, reflection,
native mutation, and dynamically loaded modules are not edge cases erased by
the set-theoretic argument above -- they are precisely the compiler soundness
obligations that Section~\ref{sec:limits} is devoted to.

\section{Implementation}
\label{sec:implementation}

The prototype consists of a Babel-based compiler, a graph/link pass, a small
runtime, and a Vite adapter. The linked compiler parses each module,
discovers state/function/component exports, resolves named imports with host
or alias resolution, computes export manifests to a fixed point, links
component placements and prop origins, and finally emits each module with its
linked facts attached. Authored ES imports remain ordinary imports as far as
the host bundler is concerned -- the linking happens entirely inside the
compiler pass, not by rewriting the module graph the bundler sees.

The implementation currently supports the following module-level reactive
chains:

\begin{itemize}
  \item imported mutable object paths and exported mutator functions;
  \item exact, receiver-bounded, and parameter-relative helper effects;
  \item expression-derived module values and exhaustive \texttt{if}/\texttt{switch}
  derivations as singleton computed entities;
  \item module-level effects with canonical imported dependencies;
  \item component definitions, prop origins, keyed-row placements, render
  callbacks, children, and selected dynamic-component candidates across files.
\end{itemize}

Opaque imported calls are marked unbounded rather than guessed at. The
current fallback commits after the known call returns normally, but it
cannot observe a mutation performed later by a callback that unknown code
retained and invokes on its own schedule -- a real limitation discussed
further below, not one the implementation currently works around. Hot
replacement removes and reinstalls module-owned index fragments and module
resources on edit.

\section{Complexity Accounting and Evaluation}
\label{sec:evaluation}

\subsection{Complexity Accounting}

Let $K$ be the canonical keys, $P$ emitted key-to-pattern entries, $L$ live
entities, $X$ materialized wildcard-pattern/entity matches, and $T_W$
entities resolved for one committed write set. The current implementation
uses $O(K+P+X)$ runtime dependency-index space, $O(L)$ entity-registry space,
and worst-case $O(L)$ dirty-set space. Once caches are warm, static
resolution costs approximately $O(|W_\sigma|+T_W)$ plus marking and
rendering; a root fallback selects $O(L)$ entities in the worst case. Live
registry changes also require maintaining the affected wildcard expansions.

This complexity accounting is deliberately more conservative than the
tempting shorthand of ``$O(0)$ static runtime memory'' set against ``$O(E)$
signals,'' and that shorthand is worth naming and rejecting explicitly rather
than letting it linger implicitly in a complexity table. Dynamic signal
systems store concrete source--observer edges; this system instead stores
static metadata and structural matches. Whether one is smaller or faster than
the other is not something either accounting can settle on paper; the
measurements below quantify it directly.

\subsection{What We Do Not Evaluate}

We do not evaluate whether \texttt{memoized-dom} outperforms complete UI
frameworks. In particular, we deliberately avoid the
\texttt{js-framework-benchmark} family of keyed-table benchmarks
\cite{js-framework-benchmark} as a headline instrument. That suite primarily
measures large-table construction, reconciliation, and framework-specific
list strategies; it conflates dependency routing with DOM diffing,
scheduling, and rendering cost in a way that does not cleanly isolate the
contribution of this paper. The measurements below instead target three
properties specific to this paper's claims: how much of a program's reactive
surface is statically resolvable, how conservative that resolution is
relative to what actually executes, and what the resulting routing mechanism
costs in isolation from rendering.

\subsection{Experimental Environment}
\label{sec:env}

All measurements were collected with Bun 1.3.14 (Node v24.3.0
compatibility) on 64-bit Windows, on an AMD PRO A12-9800B with four logical
CPUs. Routing measurements used deterministically generated topologies with
1{,}024 warmup routing operations, nine in-process timing samples per fresh
process, seven lifecycle (mount/unmount) samples per process, and three
independent fresh processes for the replicated routing result reported in
Section~\ref{sec:microbench-results}. Coverage and oracle results are
correctness measurements over deterministic inputs and do not depend on
timing. All raw results, the base commit and dirty-worktree status, and the
sampling policy are
recorded alongside the artifact; timing results in particular are reported
here as evidence from a single machine, not as a cross-platform performance
claim (Section~\ref{sec:limits}).

\subsection{Instrument 1: Corpus Coverage and Linker Ablation (RQ1, RQ4)}
\label{sec:corpus}

\subsubsection{Method}
For a corpus of programs written in the analyzed source subset of
Section~\ref{sec:source}, the compiler classifies every mutation or effectful
call site into one of the tiers of Section~\ref{sec:tiers-ref}:
exact, receiver-bounded, parameter-relative, unbounded, or rejected as
unsupported syntax. Because the compiler does not implement the hook and
component runtime semantics of frameworks such as React, arbitrary
third-party applications cannot simply be pointed at it; the corpus is a
small, purposive sample of programs written for this source subset, and we
report it as such rather than as a representative sample of real-world code
in general. The linker's specific contribution is isolated by an ablation
that disables cross-module function-summary propagation while holding the
rest of the compiler fixed -- canonical imported \emph{state} identity
remains enabled, but every imported \emph{function}'s effects are treated as
unbounded rather than resolved -- and the same classification is re-run.
This ablation reuses the same single-module analysis pass with imported
function facts forced unbounded, not a differently configured compiler entry
point, so that it isolates the linker's contribution rather than comparing
two different analyses.

\subsubsection{Results}
The corpus contains 15 accepted programs (48 modules, 274 nonblank/
noncomment lines, 88 write-effect sites) spanning seven stress categories --
aliasing and re-exports, helper-depth chains at depths one, two, and four,
branch/dynamic-key analysis stress, a multi-file derived chain, five
application-shaped case studies (todo board, inventory, permissions,
notifications, finance), a collection-mutation case, and one program that
deliberately exercises a rejected construct (a namespace import that erases
named reactive identity). Table~\ref{tab:ablation} gives the aggregate tier
distribution with and without cross-module function-summary linking.

\begin{table}[h]
\centering
\caption{Write-effect tier distribution, linked vs.\ ablated (88 sites, 15 programs)}
\label{tab:ablation}
\footnotesize
\begin{tabular}{@{}lccccc@{}}
\toprule
\textbf{Mode} & Exact & Bounded & Param-rel. & Unbounded & \textbf{Rate} \\
\midrule
Linked & 35 & 24 & 29 & 0 & \textbf{0.00\%} \\
Ablated & 19 & 8 & 21 & 40 & \textbf{45.45\%} \\
\bottomrule
\end{tabular}
\end{table}

With fixed-point cross-module function summaries enabled, \emph{zero} of the
88 accepted write-effect sites in the corpus were classified unbounded.
With cross-module function-summary propagation disabled, 40 of 88 sites
(45.45\%) degraded to the unbounded tier -- these are calls to an imported
mutator function whose effect can no longer be attributed to a canonical
state path once the summary crossing the module boundary is discarded. No
program's \emph{state}-identity resolution changed under the ablation, which
is the intended scope of this specific ablation: it isolates the
contribution of function-summary linking specifically, not of canonical
state identity, which is a separate and already-enabled mechanism in both
conditions.

\subsubsection{Paper-ready statement}
Across a purpose-built corpus of 15 accepted programs (48 modules, 274
nonblank/noncomment lines, 88 write-effect sites), fixed-point imported
function summaries reduced unbounded write-effect sites from 45.45\% to
0.00\%. This result is specific to the supported, purpose-built corpus and
is not an ecosystem-representativeness claim.

\subsection{Instrument 2: Concrete-Execution Precision Oracle (RQ2, RQ5)}
\label{sec:oracle}

\subsubsection{Method}
Static read sets are unions over syntactically reachable branches
(Section~\ref{sec:runtime}), so a selected computation can be a false
positive relative to any one concrete execution; conversely, an analysis
defect could in principle produce a false negative, which is a correctness
failure rather than a performance cost. A signal-style Proxy tracker cannot
supply ground truth for this comparison, since a Proxy intercepts property
access on an object but is structurally blind to a bare lexical read or
write of a \texttt{let}/\texttt{var} binding -- exactly the state model this
paper's source subset treats as first-class reactive state. We instead use a
second, independently implemented Babel transform that does not import or
call the compiler's own read/write analysis. It discovers module bindings,
wraps concrete lexical and object-path reads, logs actual writes (including
calls to the recognized mutating built-ins of Section~\ref{sec:source}), uses
runtime object identity to recover parameter-relative writes, and evaluates
every relevant conditional direction separately. For each executed mutation
we record $S$ (statically selected computations), $O$ (computations whose
concrete execution actually read a written key), and $C$ (computations whose
externally visible output changed), and report

\[
\text{precision} = \frac{|S \cap O|}{|S|}, \qquad
\text{recall} = \frac{|S \cap O|}{|O|}.
\]

The runner fails immediately if $O \not\subseteq S$ for any recorded case,
so that a false negative -- a concrete counterexample to
Theorem~\ref{thm:coverage}'s coverage guarantee as implemented -- cannot pass
silently. Because the oracle observes one concrete execution per run, we
report results as a \emph{dynamic-trace oracle} bounded to the executions
performed, not as complete semantic ground truth over all possible
executions; both analyses also operate over the same JavaScript syntax and
can in principle share conceptual blind spots, which is a limitation of
oracle independence discussed in Section~\ref{sec:limits}.

\subsubsection{Results}
Seven base cases were exercised across 25 concrete scenarios spanning
aliasing, four-level helper-depth propagation, branching in both directions,
independent object-path fields, cross-module derived chains, and collection
receiver mutations. Table~\ref{tab:oracle-agg} gives the aggregate result;
Table~\ref{tab:oracle-cat} breaks it down by stress category.

\begin{table}[h]
\centering
\caption{Concrete-execution oracle: aggregate result}
\label{tab:oracle-agg}
\footnotesize
\begin{tabular}{@{}lc@{}}
\toprule
Cases / scenarios & 7 / 25 \\
True positives & 36 \\
False positives & 9 \\
\textbf{False negatives} & \textbf{0} \\
\textbf{Precision} & \textbf{80.0\%} \\
\textbf{Recall} & \textbf{100.0\%} \\
\bottomrule
\end{tabular}
\end{table}

\begin{table}[h]
\centering
\caption{Concrete-execution oracle: result by stress category}
\label{tab:oracle-cat}
\footnotesize
\begin{tabular}{@{}lccccc@{}}
\toprule
\textbf{Category} & Scen. & TP & FP & Prec. & Recall \\
\midrule
Aliasing & 6 & 9 & 3 & 75.0\% & 100\% \\
Helper depth & 3 & 3 & 0 & 100\% & 100\% \\
Branching & 6 & 4 & 2 & 66.7\% & 100\% \\
Multi-file deriv. & 3 & 6 & 0 & 100\% & 100\% \\
Path precision & 4 & 8 & 4 & 66.7\% & 100\% \\
Receiver mutation & 3 & 6 & 0 & 100\% & 100\% \\
\bottomrule
\end{tabular}
\end{table}

Recall was 100\% across all 25 scenarios: every computation whose concrete
execution actually read a written key was contained in the compiler's
static selection, with zero false negatives observed. Precision was 80.0\%
overall (36 true positives, 9 false positives), and the false positives
trace to exactly two sources. First, \emph{inactive branch unions}: a static
branch union includes a location that the preceding concrete execution did
not read, so the computation is selected but its output guard
(Section~\ref{sec:runtime}) observes no change -- this accounts for the two
false positives in the branching category. Second, \emph{cross-field/module-
object over-approximation}: some aliased-object and independent-field cases
select a sibling computation in addition to the two concretely affected
computations, accounting for the remaining seven false positives, spread
across the aliasing and path-precision categories. The simpler single-binding
alias cases, four-level parameter-helper chain, derived chain, and collection
receiver mutations produced no false positives; the multiple-import-alias
cases did exhibit the cross-field over-approximation described above. No case
in any category produced a false negative.

\subsubsection{Paper-ready statement}
Across 25 concrete executions covering aliases, four-level helper
propagation, both directions of conditional dependencies, independent object
paths, collection receiver mutations, and cross-module derived chains, the
compiler achieved 100\% concrete-execution recall and 80.0\% precision (36
TP, 9 FP, 0 FN). False positives arose from inactive branch unions and
cross-field/module-object over-approximation. This is implementation
evidence from a dynamic-trace oracle, not a proof of whole-program soundness
over all executions, exceptions, opaque retained callbacks, reflection, or
dynamically loaded code.

\subsection{Instrument 3: Isolated Routing Microbenchmark (RQ3)}
\label{sec:microbench}

\subsubsection{Method}
We compare two rendering-free kernels fed the same deterministic topology
and the same synchronous propagation and path-matching rules: \emph{static
index} routing as described in Section~\ref{sec:runtime}, with exact readers
plus wildcard patterns whose live matches are materialized during lifecycle
changes; and a minimal \emph{dynamic-subscription} kernel using a
key-to-subscriber \texttt{Set} populated at mount time. The dynamic kernel is
not Solid, Vue, or Svelte and is not presented as a benchmark of either
project's actual implementation; it is a comparable reference mechanism.
Unlike PixieDust's evaluation, which compared complete PixieDust applications
against MobX/React implementations including rendering behavior
\cite{pixiedust}, this microbenchmark removes rendering entirely so that it
isolates dependency-routing cost on its own. The harness verifies identical
selected entity sets between kernels before timing. Six topologies vary key
count, entity count, fan-out, wildcard density, writes per commit, and
derived-computation depth $d_c$ (Table~\ref{tab:topologies}).

\begin{table}[h]
\centering
\caption{Generated routing topologies}
\label{tab:topologies}
\footnotesize
\begin{tabular}{@{}lcccccc@{}}
\toprule
\textbf{Topology} & Keys & Ent. & Fan & Wild.\% & W/commit & $d_c$ \\
\midrule
Exact, low fan-out & 128 & 256 & 2 & 0\% & 1 & 1 \\
Exact, high fan-out & 128 & 256 & 64 & 0\% & 1 & 1 \\
Wildcard-heavy & 128 & 256 & 32 & 75\% & 1 & 1 \\
Four writes/commit & 128 & 256 & 16 & 25\% & 4 & 1 \\
Derived depth four & 128 & 256 & 16 & 25\% & 1 & 4 \\
Large topology & 512 & 1024 & 32 & 50\% & 2 & 2 \\
\bottomrule
\end{tabular}
\end{table}

\subsubsection{Results}
\label{sec:microbench-results}
Table~\ref{tab:routing} reports the median steady-state routing time across
three fresh processes (nine timing samples each, after 1{,}024 warmup
routes). The process minimum--maximum ranges are retained in the artifact.

\begin{table}[h]
\centering
\caption{Steady-state routing: static index vs.\ dynamic subscription (median of 3 fresh processes)}
\label{tab:routing}
\footnotesize
\begin{tabular}{@{}lccc@{}}
\toprule
\textbf{Topology} & Static (ns) & Dynamic (ns) & Ratio \\
\midrule
Exact, low fan-out & 7{,}452 & 3{,}971 & 1.88$\times$ \\
Exact, high fan-out & 73{,}046 & 79{,}609 & 0.92$\times$ \\
Wildcard-heavy & 32{,}849 & 26{,}006 & 1.26$\times$ \\
Four writes/commit & 69{,}241 & 45{,}324 & 1.53$\times$ \\
Derived depth four & 42{,}693 & 31{,}518 & 1.35$\times$ \\
Large topology & 94{,}934 & 84{,}658 & 1.12$\times$ \\
\bottomrule
\end{tabular}
\end{table}

Static-index routing showed no uniform steady-state speed advantage in these
six topologies. Its median was lower for exact high fan-out (0.92$\times$)
and higher in the other five topologies (1.12$\times$--1.88$\times$) than the
minimal dynamic-subscription kernel. The picture is different at the lifecycle
boundary. Table~\ref{tab:lifecycle} reports mount cost per entity.

\begin{table}[h]
\centering
\caption{Mount cost per entity: static index vs.\ dynamic subscription}
\label{tab:lifecycle}
\footnotesize
\begin{tabular}{@{}lccl@{}}
\toprule
\textbf{Topology} & Static (ns) & Dynamic (ns) & Interpretation \\
\midrule
Exact, low fan-out & 576 & 1{,}110 & static $\approx$1.9$\times$ cheaper \\
Exact, high fan-out & 360 & 9{,}756 & static $\approx$27.1$\times$ cheaper \\
Wildcard-heavy & 127{,}541 & 3{,}722 & static $\approx$34.3$\times$ slower \\
Four writes/commit & 22{,}138 & 1{,}516 & static $\approx$14.6$\times$ slower \\
Derived depth four & 25{,}145 & 1{,}292 & static $\approx$19.5$\times$ slower \\
Large topology & 229{,}717 & 5{,}518 & static $\approx$41.6$\times$ slower \\
\bottomrule
\end{tabular}
\end{table}

Exact, high-fan-out topologies mount roughly 27$\times$ cheaper under static
routing, because precomputed exact edges avoid installing many individual
dynamic subscriber entries. Topologies containing structural wildcard
patterns reverse this advantage, because the static kernel tests and
materializes those matches at mount time -- a real
cost that a bare asymptotic accounting of ``no runtime graph'' would hide,
and exactly the cost Section~\ref{sec:runtime} names when it distinguishes
``without runtime read tracking'' from ``without runtime state.'' Logical
retained-edge counts were equal between kernels in exact topologies; in
wildcard topologies the static representation retained both the wildcard
patterns themselves and their materialized matches, adding roughly
1.5--2.3\% additional logical representation units in these topologies.

\subsubsection{Paper-ready statement}
A rendering-free mechanism benchmark found no uniform steady-state routing
advantage for the static index: across six deterministic topologies on one
evaluation machine, its median was lower in one topology (0.92$\times$) and
higher in the other five (1.12$\times$--1.88$\times$) than an equivalent
minimal dynamic-subscription kernel. Lifecycle cost depended on
representation: exact high-fan-out mounting was approximately 27$\times$
cheaper statically, while wildcard-heavy static mounting was approximately
34$\times$ slower, because live structural matches were
materialized at mount time. These results should not be used to rank
complete frameworks; they exclude DOM work, rendering, browser scheduling,
and each framework's production optimizations, and they come from a single
evaluation machine (Section~\ref{sec:env}).

\subsection{Synthesis}

Table~\ref{tab:eval-plan} maps each research question to its instrument and
headline result. The strongest defensible summary of the full evaluation is:
fixed-point ES-module linking materially improves the precision of static
reactive effect routing within the supported TypeScript/JSX subset, while
retaining every dependency observed across the tested concrete executions;
it shifts where dependency-handling cost is paid across compile, mount, and
update phases, and it is not a demonstrated inherent steady-state speedup
over dynamic subscriptions.

\begin{table}[h]
\centering
\caption{Evaluation instruments and headline results}
\label{tab:eval-plan}
\footnotesize
\begin{tabular}{@{}p{1.6cm}p{3.0cm}p{3.0cm}@{}}
\toprule
\textbf{Question} & \textbf{Instrument} & \textbf{Headline result} \\
\midrule
RQ1: coverage & Corpus classification (\S\ref{sec:corpus}) & 0.00\% unbounded, linked \\
RQ2: concrete coverage and precision & Independent oracle (\S\ref{sec:oracle}) & 100\% recall, 80.0\% precision \\
RQ3: routing cost & Static vs.\ dynamic kernel (\S\ref{sec:microbench}) & 0.92$\times$ in one case; 1.12--1.88$\times$ in five; 27$\times$ cheaper / 34$\times$ costlier mount \\
RQ4: linker value & Linker ablation (\S\ref{sec:corpus}) & 45.45\% $\rightarrow$ 0.00\% unbounded \\
RQ5: precision limits & Category breakdown (\S\ref{sec:oracle}) & branches and cross-field aliasing dominate FPs \\
\bottomrule
\end{tabular}
\end{table}

\section{Limitations and Threats to Validity}
\label{sec:limits}

\begin{itemize}
  \item \textbf{Restricted language surface.} Full JavaScript effect inference
  is undecidable in general. Dynamic keys, reflection, \texttt{eval}, native
  code, and unresolved libraries reduce precision or defeat observation.
  \item \textbf{Temporal escape.} A callback retained by opaque code can mutate
  state after the analyzed call's commit. A post-call fallback does not make
  that future write observable.
  \item \textbf{Exception semantics.} The current design does not yet specify
  and validate commit behavior for every exceptional exit.
  \item \textbf{Closed-world linking.} Dynamic imports, independently deployed
  chunks, and version skew can violate the summaries used at compilation.
  \item \textbf{Conservative branches and rows.} Static branch unions and
  wildcard row patterns can select more work than execution-sensitive signals;
  value guards suppress outputs, not all recomputation.
  \item \textbf{Scheduling.} Structural depth, early module-computed depth,
  effect phases, and drain passes are not a general static topological ordering
  proof. Cyclic behavior is detected by a pass bound.
  \item \textbf{Runtime state.} The design still needs a dependency index,
  wildcard matches, a live registry/tree, dirty state, and scheduling queues.
  \item \textbf{Mount scope.} The current connected graph admits one top-level
  mount boundary; independently scoped roots require isolated registries and
  access indexes.
  \item \textbf{Prototype and benchmark bias.} Synthetic routing topologies
  can favor either representation and do not establish application-wide
  performance or developer usability.
\end{itemize}

\section{Conclusion}

Static dependency analysis, dependency inversion, incremental browser views,
and compile-time assignment instrumentation are not new ideas -- IceDust,
PixieDust, and Svelte's own compiler have each, in their own settings,
already gotten there first. What is left to establish, and what this paper
has tried to establish carefully rather than assert, is narrower and more
specific: whether statically inferred reactive identity and imperative
effects can be linked through an ordinary, multi-file ES-module graph, and
then routed to direct DOM, component, and effect entities without paying for
dependency discovery on every read. The module-qualified model in
Section~\ref{sec:source} is built specifically to capture the cases a
single-file analysis misses -- shared state pulled into its own module,
mutation hidden behind an imported helper, derived values computed once and
consumed from several files -- and the conditional coverage theorem in
Section~\ref{sec:formal} is written the way it is so that the open questions
are visible rather than buried: analysis coverage, summary soundness, commit
boundaries, and pattern--instance consistency are named as assumptions
precisely because they are the things a future proof, or a future compiler
bug report, will have to reckon with. The artifact-backed evaluation found
that cross-module summary linking eliminated unbounded effects in the tested
corpus and retained every dependency observed by the concrete oracle, while
also exposing conservative false positives and representation-dependent
routing and lifecycle costs. Broader-corpus validation and independent
replication remain necessary before making a broader generality or performance
claim.

\begin{thebibliography}{00}

\bibitem{lustre}
N. Halbwachs, P. Caspi, P. Raymond, and D. Pilaud,
``The synchronous data flow programming language LUSTRE,''
\emph{Proceedings of the IEEE}, vol. 79, no. 9, pp. 1305--1320, 1991,
doi: 10.1109/5.97300.

\bibitem{incremental-lambda}
Y. Cai, P. G. Giarrusso, T. Rendel, and K. Ostermann,
``A theory of changes for higher-order languages: Incrementalizing
$\lambda$-calculi by static differentiation,'' in \emph{Proc. PLDI}, 2014,
pp. 145--155, doi: 10.1145/2594291.2594304.

\bibitem{icedust}
D. C. Harkes, D. M. Groenewegen, and E. Visser,
``IceDust: Incremental and eventual computation of derived values in persistent
object graphs,'' in \emph{Proc. ECOOP}, LIPIcs, vol. 56, 2016, pp. 11:1--11:26,
doi: 10.4230/LIPIcs.ECOOP.2016.11.

\bibitem{pixiedust}
N. ten Veen, D. C. Harkes, and E. Visser,
``PixieDust: Declarative incremental user interface rendering through static
dependency tracking,'' in \emph{Proc. WWW Companion}, 2018, pp. 721--729,
doi: 10.1145/3184558.3185978.

\bibitem{reactive-variables}
C. Schuster and C. Flanagan, ``Reactive programming with reactive variables,''
in \emph{Companion Proc. Modularity}, 2016, pp. 29--33,
doi: 10.1145/2892664.2892666.

\bibitem{remus}
B. Oeyen, J. De Koster, and W. De Meuter,
``Reactive programming on the bare metal: A formal model for a low-level
reactive virtual machine,'' REBLS, VUB-SOFT-TR-22-15, 2022. [Online]. Available:
\url{https://soft.vub.ac.be/Publications/2022/vub-tr-soft-22-15.pdf}

\bibitem{svelte3}
R. Harris, ``Svelte 3: Rethinking reactivity,'' 2019. [Online]. Available:
\url{https://svelte.dev/blog/svelte-3-rethinking-reactivity}

\bibitem{svelte-legacy}
Svelte, ``Reactive \texttt{\$:} statements.'' [Online]. Available:
\url{https://svelte.dev/docs/svelte/legacy-reactive-assignments}

\bibitem{svelte-runes}
The Svelte Team, ``Introducing runes,'' 2023. [Online]. Available:
\url{https://svelte.dev/blog/runes}

\bibitem{solid-reactivity}
SolidJS, ``Fine-grained reactivity.'' [Online]. Available:
\url{https://docs.solidjs.com/advanced-concepts/fine-grained-reactivity}

\bibitem{vue-reactivity}
Vue.js, ``Reactivity in depth.'' [Online]. Available:
\url{https://vuejs.org/guide/extras/reactivity-in-depth.html}

\bibitem{acar-sac}
U. A. Acar, G. E. Blelloch, and R. Harper, ``Adaptive functional programming,''
\emph{ACM TOPLAS}, vol. 28, no. 6, pp. 990--1034, 2006,
doi: 10.1145/1186632.1186634.

\bibitem{js-framework-benchmark}
S. Krause et al., ``js-framework-benchmark,'' GitHub repository. [Online].
Available: \url{https://github.com/krausest/js-framework-benchmark}

\bibitem{adapton}
M. A. Hammer, K. Y. Phang, M. Hicks, and J. S. Foster,
``Adapton: Composable, demand-driven incremental computation,'' in
\emph{Proc. PLDI}, 2014, pp. 156--166, doi: 10.1145/2594291.2594324.

\end{thebibliography}

\end{document}
